\section{Related Work}

This work sits at the intersection of Systems Biology, Applied Category Theory, and Agentic AI. While significant
research exists within each domain, the formal synthesis of biological control topologies with agentic software
architectures has received limited attention.

\subsection{Network Motifs in Systems Biology}

The concept of ``Network Motifs''---statistically over-represented sub-graphs in complex networks---was introduced
by Milo et al.~\cite{milo2002network}. Their work demonstrated that biological networks are not random but are composed of specific
building blocks selected for functional data processing. Alon~\cite{alon2007network} further characterized the dynamical properties
of these motifs, identifying the Coherent Feed-Forward Loop (CFFL) as a persistence detector. We extend this by
mapping these motifs to the stochastic nature of Generative AI.

\subsection{Applied Category Theory (ACT)}

To formalize network structure, we draw upon ACT. Fong and Spivak~\cite{fong2019seven} provide a comprehensive introduction to Applied Category Theory. Spivak~\cite{spivak2021learners} and Vagner et al.~\cite{vagner2015algebras} established a rigorous framework
for modeling Open Dynamical Systems using the category $\mathbf{Poly}$ and the Operad of Wiring Diagrams. Niu and Spivak~\cite{niu2023polynomial} develop the full mathematical theory of polynomial functors as a language for interaction. Marom et al.~\cite{marom2026category} independently relate biological and engineered stimulus--response systems by a structure-preserving functor between subcategories of dynamical systems, articulating the substrate-independence principle we adopt in \S\ref{sec:mapping}: systems sharing compositional structure can be related by interface logic alone, without sharing a physical substrate. To our
knowledge, this is the first application of Polynomial Functors specifically designed to model the interface of
LLM Agents and to verify safety properties in Agentic topologies.

\subsection{The Architecture Triple}\label{sec:archagents}

De los Riscos et al.~\cite{delosriscos2026categorical} introduce the \textsc{ArchAgents} category, whose objects
are \emph{Architecture triples} $A=(G,\mathrm{Know},\Phi)$: $G$ is the syntactic wiring (typed modules and
ports), $\mathrm{Know}$ is the structural knowledge the architecture maintains---its invariants and
\emph{certificates}---and $\Phi$ is the deployment map from abstract capability slots to concrete model and
tool implementations. A \emph{certificate} is a triple $(\tau,\sigma,\mathrm{evds})$: a theorem statement
$\tau$, a map $\sigma$ from its symbols to architecture parameters, and a mechanically replayable derivation
$\mathrm{evds}$~\cite{banu2026harness}. Morphisms are structure-preserving translations---in practice,
\emph{compilers} between frameworks. Together with the shared-interface correspondence of \S\ref{sec:mapping},
this triple is one of the two categorical objects this monograph builds on; it recurs below as the carrier of
Operon's structural guarantees.

\subsection{Reliability in Agentic AI}

Techniques such as ``Chain of Thought''~\cite{wei2022chain} utilize iterative looping to improve output quality. However, these
methods operate primarily at the level of the prompt (the input signal) rather than the topology (the wiring).
By importing the concept of Autopoiesis~\cite{maturana1980autopoiesis}, we propose a methodology where reliability is a property of the
network architecture itself.

\subsection{Epistemic Logic in Distributed Systems}

Fagin et al.~\cite{fagin1995reasoning} established the foundational framework for reasoning about knowledge in multi-agent systems, formalizing what agents know and what they know about each other's knowledge. Halpern and Moses~\cite{halpern1990knowledge} proved the impossibility of attaining common knowledge in asynchronous systems, a result with deep implications for coordination protocols. To our knowledge, these formal methods from epistemic logic have not previously been applied to the design and optimization of multi-agent AI architectures.

\subsection{Temporal Databases and Bi-Temporal Data Models}

Bi-temporal data management---tracking both \emph{valid time} (when a fact is true in the world) and \emph{transaction time} (when the system recorded it)---has been studied extensively in the database community. Snodgrass~\cite{snodgrass2000temporal} provided the foundational treatment, distinguishing valid-time, transaction-time, and bi-temporal relations, and demonstrating that the two time axes are orthogonal: a fact may be recorded before it becomes valid, corrected after it ceases to be valid, or both. SQL:2011~\cite{kulkarni2012temporal} standardized temporal query support, including \texttt{FOR SYSTEM\_TIME} and \texttt{FOR BUSINESS\_TIME} clauses.

\subsection{Adaptive Multi-Agent Assembly and Meta-Control}

Dupoux, LeCun, and Malik~\cite{dupoux2026learning} propose a three-system cognitive architecture: System~A (observational, statistical, cheap) discovers representations, System~B (action-oriented, goal-directed, expensive) discovers causal structure, and System~M (meta-control) routes data between them. This tripartite structure maps naturally onto Operon's fast/deep nucleus distinction and motivates the \texttt{WatcherComponent} as a concrete instantiation of System~M (\S\ref{sec:adaptive-structure}).

Hao et al.~\cite{hao2026bigmas} demonstrate empirically that incorrect multi-agent runs require systematically more routing decisions than successful ones (e.g., 9.4 vs 7.3 mean decisions on planning tasks). This finding grounds our use of intervention count as a convergence proxy: when the watcher's cumulative retry/escalate count exceeds a threshold relative to stage count, it emits a non-convergence signal and halts the organism (\S\ref{sec:watcher-impl}).

Jiang et al.~\cite{jiang2026anatomy} provide a structure-oriented taxonomy of Memory-Augmented Generation (MAG) systems, categorizing them into lightweight semantic, entity-centric, episodic/reflective, and structured/hierarchical designs. Their empirical analysis reveals that append-only memory architectures are significantly more robust to backbone format errors than complex structured alternatives---a finding that validates Operon's design decision to make \texttt{BiTemporalMemory} append-only. They also quantify the ``Agency Tax'' (latency and token overhead of memory maintenance), highlighting a practical concern that the convergence with operational runtimes must address.

Lin et al.~\cite{lin2026memma} propose MemMA, a multi-agent framework coordinating the full memory cycle through a planner-worker architecture with in-situ self-evolution. Their Meta-Thinker provides strategic guidance for both memory construction and retrieval, analogous to Operon's \texttt{WatcherComponent} providing intervention decisions. Their probe-based memory verification---generating synthetic questions after each session to test memory fidelity, then repairing failures immediately---complements Operon's \texttt{counterfactual\_replay()}, which detects corrections but does not actively probe for gaps.

Feng, Wang, and Zhu~\cite{feng2026selfevolving} propose Self-evolving Embodied AI, a paradigm comprising five co-evolving components: memory self-updating, task self-switching, environment self-prediction, embodiment self-adaptation, and model self-evolution. These map directly onto Operon's phased roadmap (bi-temporal memory, adaptive assembly, sleep consolidation, developmental staging, and social learning respectively), and their emphasis on multi-timescale closed-loop co-evolution aligns with Operon's per-stage (watcher), per-run (adaptive), and per-batch (consolidation) adaptation cycles.

Zhou et al.~\cite{zhou2026externalization} provide a unified review of agent externalization through four pillars---Memory, Skills, Protocols, and Harness Engineering---tracing the progression from weights-based to harness-centric agent design.  Their framework maps directly onto Operon's categorical Architecture triple: Memory corresponds to the coalgebraic state (BiTemporalMemory), Skills to objects composed via the agentic operad (SkillStage), Protocols to syntactic wiring (WiringDiagram), and the Harness to the full Architecture $(G, \mathrm{Know}, \Phi)$.  Their observation that ``agent infrastructure transforms hard cognitive burdens into manageable forms'' is precisely the claim our structural guarantee benchmarks (\S\ref{sec:empirical}) validate empirically.

Ma et al.~\cite{ma2026atomic} demonstrate that five atomic coding skills---localization, editing, unit-test generation, issue reproduction, and code review---compose without negative interference under joint reinforcement learning, an 18.7\% average improvement over task-specific optimization. This is favorable empirical evidence for the operad composition model (\S\ref{sec:operad}), with one caveat we keep explicit: the operad preserves \emph{typed and structural} properties (interface compatibility, certificate replay) by construction, but \emph{behavioral} quality under composition is empirical---our own benchmarks find mostly non-interference alongside one negative case~\cite{banu2026benchmarks}.

We apply bi-temporal data management~\cite{snodgrass2000temporal, kulkarni2012temporal} to agentic memory in \S\ref{sec:coalgebra} and \S\ref{sec:temporal-epistemics}.
